% ============================================================
% L05 — Glosarium
% ============================================================
\flowmeta{L05}{Glosarium}{Lampiran}

\begin{flowsection}{Istilah}
\begin{center}
\renewcommand{\arraystretch}{1.3}
\begin{tabular}{|p{4.2cm}|p{9.8cm}|}
\hline
\textbf{Istilah} & \textbf{Arti} \\ \hline
Storefront & Tampilan toko online untuk pelanggan. \\ \hline
Panel Admin (ERP) & Antarmuka pengelolaan toko oleh admin. \\ \hline
SKU & Kode unik identitas varian produk. \\ \hline
Varian Bertingkat & Kombinasi atribut (mis. warna+ukuran) yang menghasilkan SKU \& harga sendiri. \\ \hline
Kupon/Voucher & Kode potongan harga dengan syarat tertentu. \\ \hline
Poin Loyalitas & Poin reward yang dikumpulkan dan dapat ditukar. \\ \hline
Reservasi Stok & Penguncian stok saat pesanan dibuat agar tidak overselling. \\ \hline
Mutasi Stok & Pergerakan stok (masuk/keluar/penyesuaian). \\ \hline
GRN & Goods Receipt Note, pencatatan penerimaan barang. \\ \hline
PO & Purchase Order, pemesanan barang ke vendor. \\ \hline
HPP/COGS & Harga pokok penjualan. \\ \hline
Webhook & Panggilan otomatis dari layanan pihak ketiga ke backend. \\ \hline
Santum/Sanctum & Mekanisme token autentikasi API. \\ \hline
Role/Role Access & Peran pengguna dan hak akses menu. \\ \hline
\end{tabular}
\end{center}
\end{flowsection}

\docver{v1.0}{Sep 2026}
